\documentclass[11pt,a4paper]{article}
\usepackage[utf8]{inputenc}
\usepackage[polish]{babel}
\usepackage[T1]{fontenc}
\usepackage{graphicx}
\usepackage{hyperref}
\usepackage{listings}
\usepackage{xcolor}
\usepackage{geometry}
\usepackage{float}
\usepackage{amsmath}
\usepackage{algorithm}
\usepackage{algpseudocode}

\geometry{margin=2.5cm}

\lstset{
    basicstyle=\ttfamily\small,
    breaklines=true,
    frame=single,
    numbers=left,
    numberstyle=\tiny,
    language=Python
}

\title{Domain Logic \& Decision Engine}
\author{Fire Simulation System}

\begin{document}

\maketitle
% \tableofcontents
\newpage

\section{Wprowadzenie}

\subsection{Cel dokumentu}

Dokument ten zawiera opis logiki domenowej, algorytmów decyzyjnych, modeli symulacji oraz heurystyk wykorzystywanych w systemie.

\subsection{Zakres dokumentu}

\begin{itemize}
    \item Model fizyczny symulacji pożaru
    \item Algorytm MCTS (Monte Carlo Tree Search)
    \item Heurystyki decyzyjne
    \item Model rozprzestrzeniania się ognia
    \item System nagród i kar w MCTS
    \item Ograniczenia i założenia modelu
\end{itemize}

\section{Model symulacji pożaru}

\subsection{Przegląd}

Symulator pożaru lasu modeluje fizyczne procesy związane z pożarem, w tym:
\begin{itemize}
    \item Rozprzestrzenianie się ognia między sektorami
    \item Zmiany poziomu ognia w sektorze
    \item Proces spalania (burn level)
    \item Proces gaszenia przez brygady pożarnicze
    \item Wpływ warunków środowiskowych (wiatr, wilgotność, temperatura)
\end{itemize}

\textbf{Ważne:} Symulator działa w module \textbf{FFSim} i jest oddzielony od systemu decyzyjnego. Wszystkie decyzje podejmowane są przez system (FFBackend, FFSup), a symulator jedynie wykonuje komendy i symuluje fizykę.

\subsection{Sektor lasu}

\subsubsection{Stany sektora}

Sektor może znajdować się w jednym z trzech stanów:
\begin{itemize}
    \item \textbf{INACTIVE} -- brak ognia
    \item \textbf{ACTIVE} -- aktywny pożar
    \item \textbf{LOST} -- sektor całkowicie spalony (burn level $\geq$ 80\%)
\end{itemize}

\subsubsection{Parametry sektora}

\begin{itemize}
    \item \textbf{fireLevel} (0-100): Intensywność ognia
    \item \textbf{burnLevel} (0-100): Procent spalenia sektora
    \item \textbf{extinguishLevel} (0-100): Postęp gaszenia
    \item \textbf{temperature}: Temperatura w stopniach Celsjusza
    \item \textbf{airHumidity}: Wilgotność powietrza (0-100\%)
    \item \textbf{windSpeed}: Prędkość wiatru (m/s)
    \item \textbf{windDirection}: Kierunek wiatru (0-360 stopni)
    \item \textbf{litterMoisture}: Wilgotność ściółki (0-100\%)
    \item \textbf{co2Concentration}: Stężenie CO2 (ppm)
    \item \textbf{pm2\_5Concentration}: Stężenie PM2.5 ($\mu$g/m³)
\end{itemize}

\subsection{Aktualizacja poziomu ognia}

Poziom ognia w sektorze jest aktualizowany w każdym ticku symulacji zgodnie z formułą:

\begin{equation}
fireLevel_{t+1} = \min(fireLevel_t + \frac{fireLevel_t}{10} \cdot \alpha \cdot F - extinguishLevel_t, 100)
\end{equation}

gdzie:
\begin{itemize}
    \item $\alpha$ -- współczynnik typu sektora
    \item $F$ -- stała mnożnika poziomu ognia (domyślnie 1.0)
    \item $extinguishLevel_t$ -- poziom gaszenia w chwili $t$
\end{itemize}

Jeśli $fireLevel_{t+1} \leq 0$, sektor przechodzi w stan \textbf{INACTIVE}.

\subsection{Aktualizacja poziomu spalenia}

Poziom spalenia (burn level) jest aktualizowany zgodnie z formułą:

\begin{equation}
burnLevel_{t+1} = \min(burnLevel_t + 0.00005 \cdot fireLevel_t^3, 100)
\end{equation}

Jeśli $burnLevel_{t+1} \geq 100$, sektor przechodzi w stan \textbf{LOST} i $fireLevel$ jest ustawiany na 0.

\subsection{Poziom gaszenia}

\textbf{[Deprecated]} Poprzedni wzór
\[
extinguishLevel_t = n_{brigades} \cdot M
\]
jest nieaktualny. \textbf{[Aktualnie w kodzie]} Każda brygada w każdej pętli symulacji obniża poziom ognia sektora o stałą stawkę:
\[
fire\_level_{t+\Delta t} = \max\bigl(0,\; fire\_level_t - r \cdot \Delta t \bigr)
\]
gdzie \(r\) to \verb|_extinguishing_rate| pojedynczej brygady (\(r=5.0\) w kodzie), a \(\Delta t\) to krok czasu symulacji. 
\subsection{Wpływ ognia na parametry środowiskowe}

\subsubsection{Temperatura}

Temperatura w sektorze jest aktualizowana zgodnie z formułą:

\begin{equation}
T_{t+1} = T_t + (T_{target} - T_t) \cdot 0.5 - T_t \cdot 0.02 + \epsilon
\end{equation}

gdzie:
\begin{itemize}
    \item $T_{target} = T_{initial} + fireLevel \cdot 0.5$ -- docelowa temperatura
    \item $\epsilon \sim U(-0.5, 0.5)$ -- losowa zmiana
    \item $T_{initial}$ -- początkowa temperatura sektora
\end{itemize}

Temperatura jest ograniczona do zakresu $[10, T_{initial} + 80]$.

\subsubsection{Wilgotność powietrza}

\begin{equation}
H_{t+1} = \max(H_t - fireLevel \cdot 0.4 + \epsilon, 5)
\end{equation}

gdzie $\epsilon \sim U(-2, 2)$.

\subsubsection{Stężenie CO2}

\begin{equation}
CO2_{t+1} = \max(CO2_t + fireLevel^{1.1} - CO2_t \cdot 0.01 + \epsilon, 300)
\end{equation}

gdzie $\epsilon \sim U(-5, 5)$.

\section{Rozprzestrzenianie się ognia}

\subsection{Model rozprzestrzeniania}

Ogień rozprzestrzenia się z sektora aktywnego do sąsiednich sektorów z prawdopodobieństwem zależnym od:
\begin{itemize}
    \item Kierunku i prędkości wiatru
    \item Typu sektora docelowego
    \item Kierunku geograficznego względem wiatru
\end{itemize}

\subsection{Współczynnik rozprzestrzeniania $\beta$}

Prawdopodobieństwo zapalenia się sąsiedniego sektora obliczane jest jako:

\begin{equation}
\beta = \min(\frac{windSpeed}{40} \cdot \alpha_{target} \cdot directionCoef, 1)
\end{equation}

gdzie:
\begin{itemize}
    \item $windSpeed$ -- prędkość wiatru (m/s)
    \item $\alpha_{target}$ -- współczynnik typu sektora docelowego
    \item $directionCoef$ -- współczynnik kierunkowy 
\end{itemize}

\subsection{Współczynnik typu sektora $\alpha$}
\label{sec:alpha}

Współczynnik $\alpha$ zależy od typu sektora:
\begin{itemize}
    \item \textbf{FOREST:} $\alpha = 1.0$ (najwyższe ryzyko)
    \item \textbf{GRASS:} $\alpha = 0.8$
    \item \textbf{BUILDING:} $\alpha = 0.3$
    \item \textbf{WATER:} $\alpha = 0.0$ (niepalne)
\end{itemize}

\subsection{Współczynnik kierunkowy}
\label{sec:direction}

Współczynnik kierunkowy zależy od różnicy między kierunkiem wiatru a kierunkiem do sąsiedniego sektora:

\begin{equation}
diff = \min(|windDirection - sectorDirection|, 8 - |windDirection - sectorDirection|)
\end{equation}

Wartości współczynnika:
\begin{itemize}
    \item $diff = 0$: $directionCoef = 0.8$
    \item $diff = 1$: $directionCoef = 0.5$
    \item $diff = 2$: $directionCoef = 0.25$
    \item $diff = 3$: $directionCoef = 0.1$
    \item $diff = 4$: $directionCoef = 0.05$
    \item $diff > 4$: $directionCoef = 0.0$
\end{itemize}

\subsection{Model wiatru}

Wiatr jest modelowany jako dynamiczny parametr, który zmienia się w czasie:

\begin{itemize}
    \item \textbf{Prędkość wiatru:} $windSpeed_{t+1} = \max(0, windSpeed_t + \epsilon)$ gdzie $\epsilon \sim U(-2.5, 2.5)$
    \item \textbf{Kierunek wiatru:} zmienia się losowo o -1, 0 lub +1 jednostkę w każdym ticku (cyklicznie, 1-8)
\end{itemize}

\section{Algorytm MCTS}

\subsection{Przegląd}

Monte Carlo Tree Search (MCTS) jest wykorzystywany w module \textbf{FFSup} do generowania rekomendacji decyzyjnych. Algorytm buduje drzewo możliwych stanów systemu, symulując akcje agentów i oceniając nagrody.

\subsection{Moduł działania}

MCTS działa w module \textbf{FFSup} i wykorzystuje następujące wejścia:
\begin{itemize}
    \item Stan wszystkich sektorów (fireLevel, burnLevel, fireState)
    \item Pozycje i stany wszystkich brygad pożarniczych
    \item Pozycje i stany wszystkich patroli leśnych
    \item Konfiguracja lasu (sąsiedztwo sektorów, typy sektorów)
\end{itemize}

Wyjścia:
\begin{itemize}
    \item Lista rekomendowanych akcji: $(agentId, sectorId)$
    \item Każda rekomendacja określa, którą brygadę wysłać do którego sektora
\end{itemize}

\subsection{Faza selekcji (Selection)}

Algorytm wybiera ścieżkę od korzenia do liścia używając formuły UCT (Upper Confidence bound applied to Trees):

\begin{equation}
UCT(node, parent) = \frac{Q[node]}{N[node]} + c \cdot \sqrt{\frac{\ln(N[parent])}{N[node]}}
\end{equation}

gdzie:
\begin{itemize}
    \item $Q[node]$ -- całkowita nagroda dla węzła
    \item $N[node]$ -- liczba wizyt węzła
    \item $N[parent]$ -- liczba wizyt rodzica
    \item $c = \sqrt{2}$ -- współczynnik eksploracji (dla UCT1)
\end{itemize}

Węzeł z najwyższą wartością UCT jest wybierany do ekspansji.

\subsection{Faza ekspansji (Expansion)}

Po dotarciu do liścia, węzeł jest rozwijany przez wygenerowanie wszystkich możliwych akcji:
\begin{itemize}
    \item Dla każdej dostępnej brygady pożarniczej: możliwe akcje to przypisanie do sektora z aktywnym pożarem
    \item Dla patroli leśnych: używana jest heurystyka
\end{itemize}

\subsection{Faza symulacji (Simulation)}

Z wybranego węzła wykonywana jest losowa symulacja do stanu terminalnego (lub do maksymalnej liczby kroków). W każdym kroku:
\begin{enumerate}
    \item Wybierana jest losowa akcja dla każdego agenta
    \item Stan systemu jest aktualizowany (rozprzestrzenianie ognia, gaszenie)
    \item Sprawdzane jest, czy osiągnięto stan terminalny
\end{enumerate}

\subsection{Faza propagacji wstecznej (Backpropagation)}

Nagroda z symulacji jest propagowana wstecz do korzenia, aktualizując wartości $Q$ i $N$ dla wszystkich węzłów na ścieżce:

\begin{algorithmic}
\For{węzeł w ścieżce (od liścia do korzenia)}
    \State $N[węzeł] \gets N[węzeł] + 1$
    \State $Q[węzeł] \gets Q[węzeł] + reward$
\EndFor
\end{algorithmic}

\subsection{Strategie nagród}

System wykorzystuje trzy strategie nagród, które różnią się wagami przypisanymi do różnych komponentów:

\subsubsection{Komponenty nagród i kar}

\begin{itemize}
    \item \textbf{LOST\_SECTOR\_PENALTY} -- kara za utracone sektory (burn level $\geq$ 100\%)
    \item \textbf{EXTINGUISHED\_FIRES\_REWARD} -- nagroda za ugaszone pożary
    \item \textbf{SPREAD\_PREVENTION\_REWARD} -- nagroda za zapobieganie rozprzestrzenianiu
    \item \textbf{BURNED\_SECTOR\_PENALTY} -- kara proporcjonalna do poziomu spalenia
    \item \textbf{SECTOR\_FIRE\_LEVEL\_PENALTY} -- kara związana z intensywnością ognia
\end{itemize}

\subsubsection{Strategia agresywna}

Wysokie wagi dla:
\begin{itemize}
    \item EXTINGUISHED\_FIRES\_REWARD
    \item SPREAD\_PREVENTION\_REWARD
\end{itemize}

Niskie wagi dla kar. Skupia się na szybkim gaszeniu pożarów.

\subsubsection{Strategia zrównoważona}

Zbalansowane wagi dla wszystkich komponentów. Optymalizuje zarówno gaszenie, jak i zapobieganie rozprzestrzenianiu.

\subsubsection{Strategia konserwatywna}

Wysokie wagi dla:
\begin{itemize}
    \item SPREAD\_PREVENTION\_REWARD
    \item LOST\_SECTOR\_PENALTY
\end{itemize}

Skupia się na ochronie sektorów przed utratą.

\subsection{Obliczanie nagrody}

Całkowita nagroda dla stanu obliczana jest jako:

\begin{equation}
reward = penalties + rewards
\end{equation}

gdzie:

\begin{align}
P &=
- L \cdot \lambda_{\text{lost}}
- F \cdot \lambda_{\text{fire}}
- B \cdot \lambda_{\text{burn}}, \\[0.5em]
R &=
E \cdot \rho_{\text{ext}}
+ S \cdot \rho_{\text{spread}}.
\end{align}

gdzie:  
L — liczba utraconych sektorów, 
F — całkowita intensywność ognia, itd.

\subsection{Warunki terminalne}

Symulacja MCTS kończy się, gdy:
\begin{itemize}
    \item Wszystkie pożary zostały ugaszone (fireState = INACTIVE dla wszystkich sektorów)
    \item Wszystkie sektory zostały utracone (fireState = LOST dla wszystkich sektorów)
    \item Osiągnięto maksymalną liczbę kroków symulacji
\end{itemize}

\subsection{Ograniczenia obliczeniowe}

\begin{itemize}
    \item \textbf{Czas obliczeń:} domyślnie 1 sekunda na wywołanie MCTS
    \item \textbf{Maksymalna liczba kroków symulacji:} 100 kroków (zabezpieczenie przed nieskończoną pętlą)
    \item \textbf{Liczba workerów MCTS:} konfigurowalna (domyślnie 2)
\end{itemize}

\section{Heurystyki decyzyjne}

\subsection{Heurystyka dla patroli leśnych}
\label{sec:patrol-heuristic}

Patrole leśne używają heurystyki zamiast MCTS do wyboru sektorów do patrolowania:

\begin{enumerate}
    \item \textbf{Filtrowanie:} Ignorowane są sektory z aktywnym pożarem (fireState = ACTIVE)
    \item \textbf{Priorytetyzacja:} Priorytet mają "starsze" sektory -- te, które nie były patrolowane najdłużej
    \item \textbf{Wybór:} Każdy patrol wybiera inny sektor (brak duplikatów)
\end{enumerate}

Wiek sektora obliczany jest jako:
\begin{equation}
age = currentStep - lastPatrolStep
\end{equation}

Sektory nigdy nie patrolowane traktowane są jako najstarsze.

\subsection{Heurystyka wymaganych brygad}

Liczba brygad wymaganych do skutecznego gaszenia sektora obliczana jest jako:

\begin{equation}
requiredBrigades = \lceil \frac{fireLevel}{3} \rceil
\end{equation}

\section{Optymalizacje}

\subsection{Ograniczenie przestrzeni akcji}

Aby zmniejszyć przestrzeń akcji w MCTS, system wykorzystuje następujące optymalizacje:

\begin{itemize}
    \item \textbf{TOP\_SECTOR\_PERCENTAGE:} Rozważane są tylko sektory z najwyższym poziomem ognia (domyślnie top 10\%)
    \item \textbf{MAX\_COORDINATED\_SECTORS:} Maksymalna liczba sektorów, do których można przypisać brygady w jednym kroku (domyślnie 10)
\end{itemize}

\subsection{Klonowanie stanu}

Aby uniknąć współdzielenia stanu między węzłami MCTS, wszystkie sektory i agenci są klonowane przed użyciem w symulacji. Zapewnia to poprawność predykcji.

\section{Ograniczenia modelu}

\subsection{Założenia upraszczające}

\begin{itemize}
    \item \textbf{Model wiatru:} Uproszczony model losowych zmian, nie uwzględnia rzeczywistych danych meteorologicznych
    \item \textbf{Rozprzestrzenianie ognia:} Prawdopodobieństwo zapalenia zależy tylko od wiatru i typu sektora, nie uwzględnia topografii
    \item \textbf{Gaszenie:} Liniowa zależność między liczbą brygad a poziomem gaszenia
    \item \textbf{Agenty:} Agenty poruszają się z stałą prędkością, nie uwzględniają przeszkód terenowych
\end{itemize}

\subsection{Ograniczenia obliczeniowe}

\begin{itemize}
    \item MCTS działa w czasie rzeczywistym, więc czas obliczeń jest ograniczony
    \item Symulacja MCTS jest uproszczona w porównaniu do pełnej symulacji fizycznej
    \item Nie wszystkie możliwe akcje są rozważane (optymalizacje przestrzeni akcji)
\end{itemize}

\subsection{Charakter symulacyjny}

Wszystkie dane są generowane przez symulator:
\begin{itemize}
    \item Brak realnych czujników środowiskowych
    \item Brak realnych danych meteorologicznych
    \item Wszystkie parametry są modelowane matematycznie
\end{itemize}

\section{Architektura decyzyjna}

\subsection{Separacja logiki decyzyjnej}

Logika decyzyjna jest oddzielona od agentów. Agent reprezentuje runtime i wykonanie, podczas gdy decyzja jest wymienialnym modułem.

\subsection{Strategie decyzyjne}

System wspiera następujące strategie decyzyjne:

\begin{itemize}
    \item \textbf{Deterministyczna:} Reguły i heurystyki
    \item \textbf{MCTS:} Monte Carlo Tree Search (obecnie używane)
    \item \textbf{LLM-driven:} Wymienna strategia wykorzystująca LLM
\end{itemize}

\subsection{LLM jako strategia decyzyjna}

LLM jest traktowany jako \textbf{jedna z implementacji} modułu decyzyjnego.

\subsubsection{Decyzje LLM-eligible}

Następujące decyzje mogą być podejmowane przez LLM:
\begin{itemize}
    \item Wybór sektora do gaszenia dla brygady pożarniczej
    \item Priorytetyzacja sektorów w sytuacji wielu pożarów
\end{itemize}

\subsubsection{Decyzje deterministyczne}

Następujące decyzje pozostają deterministyczne:
\begin{itemize}
    \item Obliczanie poziomu ognia i spalenia (fizyka)
    \item Rozprzestrzenianie się ognia (model probabilistyczny)
    \item Aktualizacja pozycji agentów (fizyka ruchu)
\end{itemize}

\subsubsection{Mechanizm walidacji i fallback}

W przypadku użycia LLM:
\begin{enumerate}
    \item LLM generuje rekomendację
    \item Rekomendacja jest walidowana (sprawdzenie poprawności formatu, zakresów wartości)
    \item Rekomendacja jest logowana z informacją o źródle (LLM vs fallback)
\end{enumerate}

\section{Przepływ decyzyjny}

\subsection{Cykl decyzyjny}

\begin{enumerate}
    \item \textbf{Inicjalizacja:} System otrzymuje stan symulacji
    \item \textbf{Analiza:} FFSup analizuje stan (aktywne pożary, dostępne brygady)
    \item \textbf{Decyzja:} Wybrana strategia decyzyjna (MCTS/LLM/heurystyka) generuje rekomendacje
    \item \textbf{Walidacja:} Rekomendacje są walidowane
    \item \textbf{Publikacja:} Rekomendacje są publikowane do FFBackend
    \item \textbf{Prezentacja:} Rekomendacje są prezentowane użytkownikowi w FFVis
    \item \textbf{Wykonanie:} Użytkownik podejmuje decyzję (może zaakceptować rekomendację lub wybrać inną akcję)
    \item \textbf{Symulacja:} FFSim wykonuje akcję i aktualizuje stan
\end{enumerate}

\subsection{Moduły odpowiedzialne}

\begin{itemize}
    \item \textbf{FFSup:} Generowanie rekomendacji (MCTS, LLM, heurystyki)
    \item \textbf{FFBackend:} Agregacja i publikacja rekomendacji
    \item \textbf{FFVis:} Prezentacja rekomendacji użytkownikowi
    \item \textbf{FFSim:} Wykonanie akcji i aktualizacja stanu fizycznego
\end{itemize}

\section{Podsumowanie}

System wykorzystuje algorytmy (MCTS) i heurystyki do generowania rekomendacji decyzyjnych. Kluczowe cechy:

\begin{itemize}
    \item Separacja symulatora od systemu decyzyjnego
    \item Model fizyczny symulacji z uwzględnieniem wiatru i typu sektora
\end{itemize}

\end{document}
